\section{Methodology}
\label{sec:methodology}

This section specifies the econometric framework for forecasting cocoa futures
volatility at horizons from one day to twelve months. We first define the
volatility measure, then specify the HAR-QR model and its ENSO-augmented
extension, describe the benchmark models, and finally lay out the evaluation
framework. Table~\ref{tab:model-overview} provides an at-a-glance summary of
the five models before the detailed specifications follow.

\begin{table}[htbp]
\small
\caption{Overview of models compared in this study, ordered from simplest to
most complex.}
\label{tab:model-overview}
\noindent\hspace*{-1.25cm}\begin{tabular}{lccccc}
\toprule
 & \multicolumn{3}{c}{Benchmarks} & \multicolumn{2}{c}{Main models} \\
\cmidrule(lr){2-4}\cmidrule(lr){5-6}
Feature & Hist.\ Vol. & GARCH(1,1) & HAR-OLS & HAR-QR & HAR-X-QR \\
\midrule
Model class         & Non-param.     & ARCH-type    & Linear reg.    & Quantile reg.  & Quantile reg.  \\
HAR lags ($d/w/m$)  & No             & No           & Yes            & Yes            & Yes            \\
ENSO regressor      & No             & No           & No             & No             & Yes            \\
Quantile output     & Empirical      & Simulation   & None           & Direct         & Direct         \\
Estimation          & ---            & MLE          & OLS            & Check fn.      & Check fn.      \\
Multi-period design & Rolling window & Recursive    & Direct         & Direct         & Direct         \\
Thesis role         & Naive BM       & Param.\ BM   & HAR benchmark  & Main (RQ1)     & Main (RQ2)     \\
\bottomrule
\end{tabular}
\par\smallskip
{\footnotesize Direct = horizon-specific model estimated by direct projection;
Recursive = closed-form GARCH mean recursion, quantiles via Gaussian closed-form
(Eq.~\ref{eq:garch-quantile}); Check fn.\ = check-function (quantile) loss;
BM = benchmark; Hist.\ Vol.\ = rolling 252-day window with empirical quantiles.}
\end{table}

%------------------------------------------------------------------------------
\subsection{Volatility Measure}
\label{sec:meth-vol-measure}
%------------------------------------------------------------------------------

We measure daily volatility using the Yang-Zhang estimator, which combines
overnight and trading-session variance components:
\begin{equation}
\label{eq:yang-zhang}
\hat{\sigma}^2_{\text{YZ}} = \hat{\sigma}^2_O + k \hat{\sigma}^2_C + (1-k) \hat{\sigma}^2_{\text{RS}},
\end{equation}
where $\hat{\sigma}^2_O$ captures overnight variance, $\hat{\sigma}^2_C$
captures open-to-close variance, and $\hat{\sigma}^2_{\text{RS}}$ is the
Rogers-Satchell component that exploits the daily range. The weighting
factor $k = 0.34/(1.34 + (n+1)/(n-1))$ minimizes estimator variance.

This choice is motivated by two considerations. First, Yang-Zhang achieves
efficiency gains over close-to-close estimators and second,
it can be computed over our full 40-year daily sample, unlike realized volatility
which requires high-frequency data. We validate
the Yang-Zhang proxy against five-minute realized volatility in
Section~\ref{sec:data}, comparing rolling-window aggregates (22-day
averages) of both measures via a Mincer-Zarnowitz regression. This
comparison at the multi-day level aligns with the HAR model's use of
averaged volatility as both predictors and targets.

%------------------------------------------------------------------------------
\subsection{The HAR-QR Model}
\label{sec:meth-har-qr}
%------------------------------------------------------------------------------

Combining the HAR framework with quantile regression yields the HAR-QR
model. The HAR-QR model for horizon $h$ and quantile $\tau$ is
\begin{equation}
\label{eq:har-qr}
Q_{\tau}(\bar{\sigma}_{t:t+h} \mid \mathcal{F}_t) = \beta_0(\tau) + \beta_d(\tau) \sigma_t^{(d)} + \beta_w(\tau) \sigma_t^{(w)} + \beta_m(\tau) \sigma_t^{(m)},
\end{equation}
where $\mathcal{F}_t$ denotes the information set at time $t$. The
volatility components capture variation at three frequencies:
\begin{align}
\sigma_t^{(d)} &= \sigma_t, \label{eq:har-daily} \\
\sigma_t^{(w)} &= \frac{1}{5} \sum_{i=0}^{4} \sigma_{t-i}, \label{eq:har-weekly} \\
\sigma_t^{(m)} &= \frac{1}{22} \sum_{i=0}^{21} \sigma_{t-i}. \label{eq:har-monthly}
\end{align}

The key feature of quantile regression is that each quantile $\tau$ has its
own coefficient vector, allowing the relationship between predictors and
future volatility to differ across the distribution. For the median
($\tau = 0.50$), the coefficients may resemble those from OLS estimation.
For the upper quantile ($\tau = 0.95$), the coefficients capture how the
predictors relate to volatility during high-volatility episodes. We report three quantiles: $\tau \in \{0.50, 0.75, 0.95\}$. The median
$\tau = 0.50$ serves as an anchor to OLS-based point forecasts. The
intermediate quantile $\tau = 0.75$ traces the transition from central
to tail behaviour. The upper quantile $\tau = 0.95$ is of primary
interest, as it captures the conditional upper bound of volatility, the
regime where mean forecasts are most likely to understate actual outcomes.
Lower quantiles ($\tau = 0.05$, $0.25$) are omitted as the thesis focuses
on upper-tail risk.

The dependent variable is the average daily volatility over the forecast
horizon, following the HAR literature convention
\parencite{ghyselsDirectIteratedMultiperiod2019,
clementsHarvestingHARXVolatility2024}:
\begin{equation}
\label{eq:horizon-vol}
\bar{\sigma}_{t:t+h} = \frac{1}{h} \sum_{i=1}^{h} \sigma_{t+i}.
\end{equation}
This averaging ensures comparability across horizons and aligns with the
HAR model's construction, where weekly and monthly components are themselves
averages of daily volatility.

We adopt the \emph{direct} forecasting approach: a separate model is
estimated for each horizon $h$, using $\bar{\sigma}_{t:t+h}$ as the
dependent variable. This contrasts with \emph{iterated} forecasting, which
estimates a one-step model and iterates forward.
\textcite{ghyselsDirectIteratedMultiperiod2019} find that direct forecasting
is preferred for realized volatility models. For quantile regression,
direct estimation is particularly natural, as it produces the conditional
quantile at horizon $h$ without the distributional assumptions required to
iterate quantile forecasts. Our longest horizons (6--12 months) would
require iterating hundreds of daily steps, amplifying any model
misspecification.

We estimate HAR-QR at six horizons: $h \in \{1, 5, 22, 66, 132, 264\}$
trading days, corresponding to 1~day, 1~week, 1~month, 3~months, 6~months,
and 12~months. The short horizons (1~day to 1~month) serve to validate the
methodology against existing literature: HAR-QR has been tested at horizons
up to 22 trading days \parencite{lyocsaImprovingStockMarket2021}. The long
horizons (6 and 12 months) constitute the core contribution, testing
whether HAR-QR can produce useful forecasts at horizons where no prior study
has ventured.

At long horizons, the autoregressive signal from lagged volatility fades.
We therefore augment the HAR-QR model with the ENSO climate index:
\begin{equation}
\label{eq:har-x-qr}
Q_{\tau}(\bar{\sigma}_{t:t+h} \mid \mathcal{F}_t) = \beta_0(\tau) + \beta_d(\tau) \sigma_t^{(d)} + \beta_w(\tau) \sigma_t^{(w)} + \beta_m(\tau) \sigma_t^{(m)} + \gamma(\tau) \text{ENSO}_t,
\end{equation}
We refer to this augmented specification as the \emph{HAR-X-QR} model,
where the ``-X'' suffix denotes the inclusion of the exogenous ENSO regressor.
Here $\text{ENSO}_t$ is the Ni\~no~3.4 sea surface temperature anomaly
index from NOAA, measured monthly. The coefficient $\gamma(\tau)$ captures
the marginal effect of ENSO conditions on future volatility at quantile
$\tau$.

We use \emph{contemporary} ENSO, meaning $\text{ENSO}_t$ predicts
$\bar{\sigma}_{t:t+h}$, rather than explicitly lagged ENSO values. This
specification is standard in the climate-commodity literature
\parencite{bouriNinoForecastabilityOilprice2021,
bonatoNinoNinaForecastability2023} and is causal: the Ni\~no~3.4
anomaly for month $t$ is publicly observable before any forecast is issued,
so no forward-looking information enters the model. The sensitivity of this
specification to alternative ENSO lags is examined in
Section~\ref{sec:res-robustness}.

Positive Ni\~no~3.4 anomalies indicate El~Ni\~no conditions; negative
anomalies indicate La~Ni\~na. Full details of the index construction,
sample coverage, and monthly-to-daily alignment are provided in
Section~\ref{sec:data-enso}.

The HAR-QR model is estimated by minimizing the check function loss:
\begin{equation}
\label{eq:har-qr-estimation}
\hat{\beta}(\tau) = \argmin_{\beta} \sum_{t=1}^{T-h} \rho_{\tau}(\bar{\sigma}_{t:t+h} - X_t' \beta),
\end{equation}
where $X_t = (1, \sigma_t^{(d)}, \sigma_t^{(w)}, \sigma_t^{(m)})'$ for the
basic model, or
$X_t = (1, \sigma_t^{(d)}, \sigma_t^{(w)}, \sigma_t^{(m)}, \text{ENSO}_t)'$
for the extended model. The check function is defined as:
\begin{equation}
\label{eq:check-function}
\rho_{\tau}(u) = u \cdot (\tau - \mathbf{1}_{u < 0}) =
\begin{cases}
\tau \cdot u & \text{if } u \geq 0, \\
(\tau - 1) \cdot u & \text{if } u < 0.
\end{cases}
\end{equation}
This asymmetric loss function penalizes under-prediction more heavily when
$\tau$ is large and over-prediction more heavily when $\tau$ is small,
ensuring that the estimated quantile achieves the target coverage rate. The
minimization is a linear programming problem solved efficiently using the
simplex algorithm \parencite{koenker1987algorithm}. Standard errors are
computed using the bootstrap.

%------------------------------------------------------------------------------
\subsection{Benchmark Models}
\label{sec:meth-benchmarks}
%------------------------------------------------------------------------------

We compare HAR-QR and HAR-X-QR against three benchmarks that span the major modeling
paradigms discussed in the literature review.

The first benchmark is \emph{rolling historical volatility} with a one-year
window ($w = 252$):
\begin{equation}
\label{eq:hist-vol}
\hat{\sigma}_{t+h}^{\text{hist}} = \sqrt{\frac{1}{w} \sum_{i=1}^{w} (r_{t-i+1} - \bar{r})^2}.
\end{equation}
For quantile forecasts, we compute empirical quantiles of the historical
volatility distribution: sort the $w$ past volatilities and select the
value at position $\lceil \tau \cdot w \rceil$. This model captures only
unconditional variation and serves as a naive baseline.

The second benchmark is the \emph{GARCH(1,1)} model
\parencite{bollerslev1986generalized}:
\begin{equation}
\label{eq:garch}
\sigma_t^2 = \omega + \alpha \varepsilon_{t-1}^2 + \beta \sigma_{t-1}^2.
\end{equation}
For mean forecasts at horizon $h$, the GARCH(1,1) produces closed-form
predictions:
\begin{equation}
\label{eq:garch-forecast}
\E_t[\sigma_{t+h}^2] = \bar{\sigma}^2 + (\alpha + \beta)^{h-1}(\sigma_{t+1}^2 - \bar{\sigma}^2),
\end{equation}
where $\bar{\sigma}^2 = \omega / (1 - \alpha - \beta)$ is the unconditional
variance and $(\alpha + \beta)^{h-1}$ governs the speed of mean reversion.
At long horizons, all $h$-step forecasts converge to $\bar{\sigma}^2$,
irrespective of current conditions, a structural limitation of daily
GARCH at 6--12 month horizons.

For quantile forecasts we adopt a closed-form Gaussian approach consistent
with the normal-innovation assumption used in estimation.
The mean forecast $\hat{\sigma}_{t,h}$ is obtained as the square root of
the average $h$-step variance, $\hat{\sigma}_{t,h} = \sqrt{h^{-1}\sum_{k=1}^{h}\E_t[\sigma_{t+k}^2]}$.
Forecast uncertainty is captured by the in-sample standard deviation
$s_t$ of daily volatility over the estimation window. The $\tau$-quantile
forecast is then
\begin{equation}
\label{eq:garch-quantile}
\hat{Q}_{\tau,t}^{\text{GARCH}} = \hat{\sigma}_{t,h} + \Phi^{-1}(\tau)\, s_t,
\end{equation}
where $\Phi^{-1}(\tau)$ is the standard normal quantile function.

The third benchmark is \emph{HAR-OLS}, the standard HAR model estimated
by ordinary least squares. OLS yields only a conditional mean forecast
$\hat{\mu}_{t,h}$, but a distributional forecast can be constructed from
the in-sample residuals. Let $\hat{\varepsilon}_i = \sigma_i - \hat{\mu}_i$
denote the residuals from the training window. The $\tau$-quantile forecast
is then
\begin{equation}
\label{eq:har-ols-quantile}
\hat{Q}_{\tau,t}^{\text{OLS}} = \hat{\mu}_{t,h} + \hat{F}^{-1}_{\varepsilon}(\tau),
\end{equation}
where $\hat{F}^{-1}_{\varepsilon}(\tau)$ is the empirical $\tau$-quantile of
the training-window residuals. This pseudo-quantile approach requires no
distributional assumption and produces forecasts that are directly comparable
to HAR-QR: if HAR-QR provides no improvement over HAR-OLS at quantile $\tau$,
then conditioning on the HAR structure suffices and quantile regression adds
no value. Comparing the two at $\tau = 0.50$ tests symmetry of the conditional
distribution; comparing at $\tau = 0.95$ tests whether the upper tail is
better captured by explicit quantile estimation.

%------------------------------------------------------------------------------
\subsection{Evaluation Framework}
\label{sec:meth-evaluation}
%------------------------------------------------------------------------------

Evaluating volatility forecasts is complicated by the fact that volatility
is unobserved, and every evaluation compares forecasts against an imperfect
proxy \parencite{pattonVolatilityForecastComparison}. We evaluate forecasts
against Yang-Zhang volatility computed from realized prices, using loss
functions that \textcite{pattonVolatilityForecastComparison} has shown to
be robust to noise in the volatility proxy.

For mean and median forecasts, we use QLIKE as the primary metric. All five models are therefore compared on QLIKE: the three benchmarks contribute their mean forecasts, while HAR-QR and HAR-X-QR contribute their median ($\tau = 0.50$) forecasts, which provide a natural reference point for comparison with OLS-based point forecasts.
\begin{equation}
\label{eq:qlike}
\text{QLIKE} = \frac{1}{T} \sum_{t=1}^{T} \left( \log \hat{\sigma}_t^2 + \frac{\sigma_t^2}{\hat{\sigma}_t^2} \right).
\end{equation}
\textcite{pattonVolatilityForecastComparison} derives necessary and
sufficient conditions for the loss function to yield forecast rankings that
are robust to proxy noise, showing that among nine commonly used loss
functions, only MSE and QLIKE satisfy these conditions. We prefer QLIKE
because it penalizes under-prediction of volatility proportionally more than
MSE, which is appropriate when under-prediction is costlier than
over-prediction, a property particularly relevant for upper-quantile
forecasts. As \textcite{pukeCoherentForecastingRealized2026} demonstrate,
aligning estimation and evaluation loss functions, what they call
\emph{coherent forecasting}, yields significant accuracy gains. Our HAR-QR
models are coherent by construction: the check function
(Equation~\ref{eq:check-function}) serves as both the estimation objective
and the evaluation metric.

As a robustness check, we additionally report mean squared error (MSE) and
mean absolute error (MAE) for all point-forecast models:
\begin{align}
\text{MSE} &= \frac{1}{T} \sum_{t=1}^{T} \bigl(\sigma_t - \hat{\sigma}_t\bigr)^2,
\label{eq:mse} \\
\text{MAE} &= \frac{1}{T} \sum_{t=1}^{T} \bigl|\sigma_t - \hat{\sigma}_t\bigr|.
\label{eq:mae}
\end{align}
MSE penalizes large errors quadratically and is a member of the
proxy-robust class alongside QLIKE \parencite{pattonVolatilityForecastComparison};
including it confirms that QLIKE rankings are not sensitive to the choice of
loss function within that class. MAE penalizes all deviations linearly and is
not proxy-robust, but provides an interpretable measure of average forecast
error in the original volatility units.

For quantile forecasts, we use quantile loss:
\begin{equation}
\label{eq:quantile-loss}
\text{QL}_{\tau} = \frac{1}{T} \sum_{t=1}^{T} \rho_{\tau}(\sigma_t - \hat{Q}_{\tau,t}).
\end{equation}

As a complementary diagnostic, we assess the calibration of quantile
forecasts through their \emph{empirical violation rate}. For a given
quantile level~$\tau$, the violation rate measures the fraction of
out-of-sample observations that exceed the quantile forecast:
\begin{equation}
\label{eq:violation-rate}
\widehat{VR}_{\tau} = \frac{1}{T}\sum_{t=1}^{T} \mathbf{1}\!\left\{\sigma_t > \hat{Q}_{\tau,t}\right\},
\end{equation}
where $\mathbf{1}\{\cdot\}$ is the indicator function. A well-calibrated
$\tau$-quantile forecast should satisfy $\widehat{VR}_{\tau} \approx 1 - \tau$.
For example, the $\tau = 0.95$ quantile should be exceeded approximately
$5\%$ of the time; systematic departures indicate miscalibration. Violation
rates above $1 - \tau$ signal that the model underestimates the quantile
(too narrow), while rates below $1 - \tau$ indicate overestimation
(too conservative). This diagnostic complements quantile loss, which
measures average forecast accuracy but does not directly reveal whether
the nominal coverage level is achieved.

To determine whether differences in forecast accuracy are statistically
significant, we use the Model Confidence Set (MCS) of
\textcite{hansen2011model}. The MCS identifies the set of models
$\mathcal{M}^*$ that contains the best model with a given confidence level,
avoiding the need to specify a single benchmark. Starting from the full
collection of competing models, the procedure iteratively tests for equal
predictive ability and eliminates the worst-performing model until
non-rejection occurs. We implement MCS with 10,000 stationary bootstrap
replications using the max-$t$ statistic; models in $\mathcal{M}^*$ at
the 90\% confidence level are considered statistically equivalent.

The forecast evaluation is conducted out-of-sample using a rolling estimation
window of 2,520 trading
days (approximately 10 years). At each forecast origin $t$, we estimate
parameters using observations from $t - 2519$ to $t$, generate forecasts,
and roll forward by one day. The window rolls \emph{daily} to maximize the number of
forecast--realization pairs for statistical evaluation. With approximately
30 years of out-of-sample data, this yields thousands of pairs for
statistical testing.

Several assumptions underpin the methodology. First, we assume \emph{local
stationarity}: volatility dynamics are stable within each estimation window,
though they may shift across windows. The rolling-window design
accommodates gradual parameter changes. Second, we assume that the
Yang-Zhang estimator provides a reasonable proxy for latent volatility,
which we validate empirically in Section~\ref{sec:data}. Third, we maintain
\emph{linear conditional quantiles} for parsimony, acknowledging that
non-linear extensions could potentially improve forecasts. Fourth, we
assume ENSO is exogenous to the cocoa market, which is physically
reasonable.

Based on the literature, we expect that at short horizons (1~day to
1~month), HAR-QR will match or outperform benchmarks. At long horizons
(6--12~months), standard HAR-QR will lose predictive power; we test whether
HAR-X-QR augmented with ENSO recovers predictive value, examining whether
climate information provides incremental value precisely where the
autoregressive signal fades.

Should HAR-X-QR outperform HAR-QR at long horizons, a single full-sample
comparison cannot by itself establish whether the gain reflects a structural
climate channel or a sample-specific coincidence. We therefore subject any
ENSO improvement to three diagnostics, reported in
Section~\ref{sec:res-robustness}. First, a \emph{temporal-stability} check
traces the cumulative ENSO gain as the evaluation horizon advances through
time using an expanding-window diagnostic. For each calendar month~$T$ in
the out-of-sample period, the quantile loss difference
$\Delta\text{QL}(T) = \text{QL}_{\text{HAR-X-QR}}(\{t \leq T\}) -
\text{QL}_{\text{HAR-QR}}(\{t \leq T\})$ is computed over all
forecast--realisation pairs up to and including~$T$. A structural ENSO--volatility channel should produce a
steadily accumulating gain throughout the evaluation period, not a gain
concentrated in a single episode. Second, a \emph{lag-structure}
check re-estimates HAR-X-QR with the Ni\~no~3.4 index lagged 0, 1, 3, and~6
months. Under the agronomic-transmission hypothesis the predictive content should
strengthen at lags of 6--12 months, matching the window over which
ENSO-driven rainfall anomalies are hypothesised to affect harvests. This
is consistent with the broader finding in the literature that ENSO's
predictive value for commodity volatility concentrates at longer horizons:
\textcite{bouriNinoForecastabilityOilprice2021} document gains at 2--4 year
horizons for oil-price realized volatility, while
\textcite{bonatoNinoNinaForecastability2023} find a general tendency for
ENSO predictability to strengthen at longer forecast horizons across 16
agricultural commodities. Third, a \emph{cross-commodity} check applies
the identical HAR-X-QR specification to Arabica coffee and raw sugar, two
tropical soft commodities that share weather exposure and downstream markets
with cocoa \parencite{geVolatilityMugCup2025}, since a genuine climate
channel should produce comparable gains across commodities grown in the same
equatorial belt
\parencite{ubilavaRoleNinoSouthern2018, bonatoNinoNinaForecastability2023}.
Because continuous futures price series for coffee and sugar are available
from approximately 2000 onward, the out-of-sample evaluation windows for these
commodities span roughly 15 years --- approximately half the cocoa sample.
The cross-commodity results should be interpreted in light of this shorter
evaluation period.
